\documentclass[11pt]{article}
\usepackage[margin=1in]{geometry}

\usepackage{amsmath,amssymb}
\usepackage{array}
\usepackage{longtable}
\usepackage{booktabs}
\usepackage{hyperref}

\setlength{\parindent}{0pt}
\setlength{\parskip}{1\baselineskip}

\newcommand{\code}[1]{\texttt{\detokenize{#1}}}
\newcommand{\breakcode}[1]{\begingroup\urlstyle{tt}\nolinkurl{#1}\endgroup}

\title{Gate12C-1:\\
A Higher-Order Negative Control for Compressed-Overlap\\
Parenthesization Sensitivity in LLM Replay Artifacts}
\author{Aoi Kawasaki}
\date{Draft skeleton, July 2026}

\begin{document}
\maketitle

\begin{abstract}
Gate12C-1 is a predeclared higher-order negative-control extension over
existing Gate12A replay artifacts. It tests whether compressed-overlap
parenthesization sensitivity exceeds a matched spectrum-preserving orientation
null under the frozen equal-rank protocol. The canonical twelve-source-run
surface completed with execution status \code{complete} and grid outcome
\code{no_directional_support}: all 24 run/q endpoints were informative and
coverage-complete, and all \code{q_support} and \code{run_support} calls were
false. The result closes the associator-excess hypothesis under this frozen
protocol. It is not a physical-structure claim, not a model-quality result, not
a correctness classifier, and not positive evidence for Gate12B.
\end{abstract}

\section{Introduction}
\label{sec:introduction}

Structural audits of LLM replay artifacts can surface relation-patterned
signals that are hidden by final-output scalar scores. They can also overfit
attractive structural stories if every new probe is treated as support. This
technical note records Gate12C-1 as a guardrail against that failure mode.

Gate12C-1 was designed after the Gate12B artifact audit line as a stricter
higher-order test. It asks whether the real Gate12A artifact cycles are more
parenthesization-sensitive than a matched orientation-randomized,
edge-spectrum-preserving null. The answer under the frozen equal-rank protocol
was no.

The contribution is therefore narrow. Gate12C-1 does not report a new positive
interpretability discovery. It records a completed null-controlled test that
rejects a tempting higher-order associator-excess reading under the declared
protocol. The value is claim discipline: the same audit stack that can report a
bounded Gate12B closure-signature result can also decline a stronger claim when
the matched-null test does not support it.

\section{Background}
\label{sec:background}

\subsection{Gate12A Source Artifacts}

Gate12C-1 operates on already materialized Gate12A replay artifacts. It does not
rerun model inference, modify model weights, change source text, or create a new
model-evaluation benchmark. The artifact surface consists of local objects,
directed transport relations, triangle cycles, and derived metadata already
present in the Gate12A line.

\subsection{Gate12B and the Positive Paper Line}

Gate12B remains a separate bounded positive artifact-level result. It reports an
archive-family observer-relative closure-signature pattern in existing replay
artifacts. Gate12C-1 does not add empirical support to Gate12B, does not expand
Gate12B's empirical scope, and does not invalidate the Gate12B
closure-signature claim.

\subsection{Gate12C-0 Feasibility and Gate12C-1 Test}

Gate12C-0 established that an equal-rank associator surface was operationally
feasible over the canonical real-artifact grid. Gate12C-1 then tested a stricter
higher-order hypothesis: whether compressed-overlap parenthesization sensitivity
exceeded a matched spectrum-preserving orientation null under a predeclared
hierarchy.

\section{Method}
\label{sec:method}

\subsection{Compressed-Overlap Parenthesization Probe}

Gate12C-1 reads local compressed-overlap transports derived from the existing
artifact surface. The probe compares parenthesization-sensitive compositions of
the transport objects and records whether the observed surface exceeds the
matched null according to the frozen support rule.

This note intentionally keeps the method operational. It does not introduce new
theory vocabulary or additional equations beyond the tracked protocol records.
The authoritative protocol and result memory are the workstream records listed
in Section~\ref{sec:reproducibility}.

\subsection{Equal-Rank Surface}

The Gate12C-1 alpha protocol uses an equal-rank construction. The result should
not be read as a rectangular rank-mismatch result, a statement about all possible
local-object constructions, or a statement about every possible null family.

\subsection{Matched Orientation Null}

The null preserves each edge's singular spectrum while randomizing orientation.
The primary question is whether the real artifact cycles exceed this matched
orientation-randomized, edge-spectrum-preserving null. This is a directional
excess test, not a test of positive associator magnitude alone.

\subsection{Hierarchy and Support Rule}

The frozen hierarchy reports endpoint-level and run-level support calls. Under
the recorded execution, every endpoint was informative and coverage-complete,
and all support calls were false.

\section{Result}
\label{sec:result}

The Gate12C-1 run completed over the canonical twelve Gate12A source-run
surface. Table~\ref{tab:primary-result} records the primary result fields used
by this paper skeleton.

\begin{table}[ht]
\centering
\small
\caption{Gate12C-1 primary result boundary.}
\label{tab:primary-result}
\begin{tabular}{|>{\raggedright\arraybackslash}p{0.45\textwidth}|>{\raggedright\arraybackslash}p{0.42\textwidth}|}
\hline
Field & Value \\
\hline
Execution status & \code{complete} \\
\hline
Grid outcome & \code{no_directional_support} \\
\hline
Supporting run count & 0 \\
\hline
Q-discordant run count & 0 \\
\hline
Coverage-limited endpoint count & 0 \\
\hline
Endpoint family & 24 run/q endpoints \\
\hline
Coverage & 24/24 informative and coverage-complete \\
\hline
Support calls & All \code{q_support} and \code{run_support} calls false \\
\hline
\end{tabular}
\end{table}

The clean interpretation is that the real Gate12A artifact cycles are not more
parenthesization-sensitive than the matched orientation-randomized,
edge-spectrum-preserving null under the frozen equal-rank protocol.

\section{Interpretation}
\label{sec:interpretation}

Gate12C-1 is a negative-control result. It prevents a stronger
associator-excess hypothesis from being reported as evidence when the
predeclared matched-null test did not support it.

This interpretation is deliberately bounded:

\begin{itemize}
    \item positive associator magnitude alone is not evidence;
    \item Gate12C-1 does not provide a higher-order structure discovery;
    \item Gate12C-1 does not support a physical-structure claim;
    \item Gate12C-1 is not positive evidence for Gate12B;
    \item Gate12C-1 is not an invalidation of the Gate12B result.
\end{itemize}

The result closes the specific directional-excess reading tested by the frozen
equal-rank protocol. It does not close every possible associator construction,
null choice, local-object construction, model family, or future prospective
suppression/coherence hypothesis.

\section{Threats to Validity}
\label{sec:threats}

The primary limitations are:

\begin{itemize}
    \item \textbf{Current source-run surface only.} The result is bounded to the
    canonical twelve Gate12A source-run surface used by the frozen plan.
    \item \textbf{One null family.} The primary null is the matched
    spectrum-preserving orientation null. Other nulls require separate
    predeclared plans.
    \item \textbf{Equal-rank protocol only.} The result does not report a
    rectangular rank-mismatch associator.
    \item \textbf{No source-facing overlay.} No Gate12B overlay is consumed here.
    Any overlay would be descriptive-only and would require a separate plan.
    \item \textbf{No model evaluation claim.} The result is not a model-quality
    ranking, safety proof, answer-quality label, or correctness classifier.
    \item \textbf{No weight-level mechanism.} The paper skeleton makes no causal
    claim about model weights or internal mechanisms.
\end{itemize}

\section{Relation to Gate12B}
\label{sec:relation-gate12b}

Gate12B and Gate12C-1 test different hypotheses. Gate12B reports a bounded
observer-relative closure-signature result over existing artifacts. Gate12C-1
tests a stricter higher-order compressed-overlap parenthesization-excess
hypothesis above a matched orientation null.

The negative Gate12C-1 result narrows claim discipline around the Gate12B paper
line, but it does not act as Gate12B evidence. It also does not erase the
Gate12B result. The two lines should remain connected only by explicit boundary
language or by a later predeclared descriptive compatibility plan.

\section{Reproducibility Boundary}
\label{sec:reproducibility}

This skeleton records paper-source wording only. It does not include generated
run artifacts, refresh a Zenodo release package, update package integrity files,
or rerun any Gate12A, Gate12B, or Gate12C tool.

Tracked source records:

\begin{itemize}
    \item \breakcode{workstream/236_GATE12C1_FIRST_EMPIRICAL_RESULT_MEMO.md}
    \item \breakcode{workstream/237_GATE12C1_RESULT_CLOSEOUT_AND_NEXT_BOUNDARY.md}
    \item \breakcode{workstream/238_GATE12C1_NEGATIVE_CONTROL_AND_PUBLIC_POSITIONING_PLAN.md}
\end{itemize}

Existing \breakcode{zenodo-release-*} directories remain frozen historical
release snapshots. This skeleton is not part of a refreshed Zenodo package and
does not update any \breakcode{CHECKSUMS-SHA256.txt} file.

\section{Non-Claims}
\label{sec:nonclaims}

This note explicitly does not claim:

\begin{itemize}
    \item a physical-structure discovery;
    \item a safety or deployment guarantee;
    \item a model-quality ranking;
    \item a correctness classifier;
    \item a weight-level causal mechanism;
    \item Gate12B overlay evidence;
    \item a rectangular rank-mismatch result.
\end{itemize}

\section{Conclusion}
\label{sec:conclusion}

Gate12C-1 closes the associator-excess hypothesis as negative under the frozen
equal-rank protocol. The completed result is
\code{no_directional_support}, with all 24 run/q endpoints informative and
coverage-complete and all support calls false. The result is useful because it
shows that the audit stack can reject an attractive higher-order extension
rather than automatically convert it into a positive claim. Any next empirical
work requires a new predeclared plan.

\end{document}
